\section{Boundary cases and the role of the hypotheses}\label{app:boundaries}
\begin{proposition}[Three noninterchangeable ensembles]\label{prop:boundaries}
An invariant ensemble of a frozen table, the actual ensemble at a fixed call, and the preparation-injected terminal ensemble need not coincide. Rowwise static quantization is also insufficient for global finite-state calibration.
\end{proposition}
\begin{proof}
For the identity transition on two labels, both point masses are invariant. Starting at the first label never realizes the second invariant law; its injected geometric resolvent remains the first point mass for every $s>0$. For a deterministic two-cycle started at the first label, the geometric terminal weights after a first call are $1/(1+\beta)$ and $\beta/(1+\beta)$, whereas the invariant law has equal weights and fixed-call laws alternate. These elementary examples do not satisfy the smoothing theorem's density assumption; they isolate the invalid initialization inference that the general construction avoids.

For the pooling assertion, take two equally weighted incoming rows, each supported on one scalar prediction, $0$ or $1$. With one output in each row their static distortions are zero. Sending both rows to the same global label gives conditional mean $1/2$ and squared error $1/4$. Local row minimizers do not provide a compatible global program. In the present proof the fixed-point equations hold jointly for all rows and the initial injection.
\end{proof}

A stopping rule that selects reports changes their conditional law. For instance, stop on the first positive report in an independent symmetric sequence; the stopping report is positive with probability one, not distributed as an unconditional row. This explains the conditional independence requirement on the hazard in every all-program lower. The capped independent hazard retains a deterministic-time mixture; arbitrary observation-dependent stopping does not. At $s=1$ the construction has one raw acquisition and the fixed-point map is constant, so no singular $\beta^{-1}$ convention is required. At $s\downarrow0$ the statistical constants remain uniform under raw overlap, while numerical tolerances in \eqref{eq:numerical} must generally become finer. This is the separation between acquired geometry and complete bit/time complexity.


\begin{proposition}[Quantitative separation of two stability certificates]\label{prop:stability}
Let (C) mean a one-step likelihood curvature bound $\tfrac12\Phi''\le b<1$ in a scalar predictive coordinate, and (P) mean a strict same-report contraction in a metric bi-Lipschitz to that coordinate on the whole certified domain. Neither certificate implies the other, even for positive instruments with state-dependent report laws.
\end{proposition}
\begin{proof}
For (P) without (C), use the binary Gaussian instrument of Section~\ref{sec:benchmark}, now with transition off-diagonal $a=10^{-4}$ and multiplier $\lambda=1-2a$. Bayes updating adds $2y$ to log odds, so is an isometry in the scalar Hilbert metric. Direct differentiation of the log-odds action of the transition bounds its derivative by $\lambda<1$. On the compact interval $[a,1-a]$ this metric is bi-Lipschitz to the scalar coordinate, proving (P).

At incoming probability $p=a$, the curvature coefficient is
\[
 b(a)=\lambda^2\int\frac{f_+(y)^2f_-(y)^2}{[a f_+(y)+(1-a)f_-(y)]^3}\,dy.
\]
Restrict $y$ to $[0,1]$. Since $f_+/f_-=e^{2y}$ and $f_+^2/f_-=e^4\phi(y-3)$,
\[
 b(a)\ge\frac{\lambda^2e^4}{(1+ae^2)^3}\int_2^3\phi(t)\,dt>1.
\]
For an explicit elementary verification, use $\lambda^2>998/1000$, $e^4>54$, $e^2<8$, and the decreasing-density lower sum
$\int_2^3\phi\ge16^{-1}\sum_{j=1}^{16}\phi(2+j/16)>19/1000$.
The last strict bound follows by $\sqrt{2\pi}<251/100$ and bounding each positive exponential series above by its first 31 terms plus its geometric tail. These rational inequalities give a lower greater than $1.02$. Thus (C) fails on the stated domain, without invalidating (P).

For (C) without (P), let $u\in[0,1]$, $r=1/8$, and give a two-state instrument a report $+$ with likelihood $r(1+u)$ and posterior $F_+(u)=2u/(1+u)$. It is the positive diagonal instrument with weights $r,2r$. The other report has the complementary likelihood and resets the physical preparation to probability $1/2$. Both report likelihoods are strictly positive. The integrated squared-response potential is
\[
 \Phi(u)=\frac{4ru^2}{1+u}+\frac{1-r(1+u)}4,
 \qquad \tfrac12\Phi''(u)=\frac{4r}{(1+u)^3}\le\tfrac12.
\]
Hence (C) holds. But $F_+$ fixes zero and has derivative two there. More explicitly,
$F_+^k(u)=2^ku/[1+(2^k-1)u]$.
If a metric bounded above and below by positive multiples of $|u-v|$ contracted every report by a factor $\kappa<1$, compare zero with $u=2^{-k}$ after $k$ successive positive reports. The final Euclidean distance exceeds $1/2$, while the contracted metric distance is bounded by a constant times $(\kappa/2)^k$, a contradiction. The positive-report word has positive probability and cannot be excluded. This latter atomic experiment isolates the stability issue; it is not claimed to satisfy the full-dimensional acquired-density hypothesis of the main theorem.
\end{proof}
